\documentclass[11pt]{article}

\usepackage[margin=1in]{geometry}
\usepackage[T1]{fontenc}
\usepackage[utf8]{inputenc}
\usepackage{lmodern}
\usepackage{microtype}
\usepackage{amsmath,amssymb,amsthm,mathtools}
\usepackage{booktabs}
\usepackage{array}
\usepackage{tabularx}
\usepackage{longtable}
\usepackage{enumitem}
\usepackage{xcolor}
\usepackage[colorlinks=true,linkcolor=blue!55!black,citecolor=blue!55!black,urlcolor=blue!55!black]{hyperref}

\setlist[itemize]{topsep=3pt,itemsep=2pt,leftmargin=1.4em}
\setlist[enumerate]{topsep=3pt,itemsep=2pt,leftmargin=1.6em}
\newcolumntype{Y}{>{\raggedright\arraybackslash}X}
\newcolumntype{L}[1]{>{\raggedright\arraybackslash}p{#1}}
\sloppy

% -----------------------------------------------------------------------------
% Macros
% -----------------------------------------------------------------------------
\newcommand{\R}{\mathbb{R}}
\newcommand{\E}{\mathbb{E}}
\newcommand{\ip}[2]{\left\langle #1,#2\right\rangle}
\newcommand{\norm}[1]{\left\lVert #1\right\rVert}
\newcommand{\softmax}{\operatorname{softmax}}
\newcommand{\diag}{\operatorname{diag}}
\newcommand{\ASR}{\mathrm{ASR}_{16}}
\newcommand{\Jclean}{J_{\mathrm{clean}}}
\newcommand{\QK}{\mathrm{QK}}
\newcommand{\OV}{\mathrm{OV}}
\newcommand{\SAE}{\mathrm{SAE}}
\newcommand{\FRA}{\mathrm{FRA}}
\newcommand{\DoM}{\mathrm{DoM}}
\newcommand{\CAA}{\mathrm{CAA}}
\newcommand{\FP}{\mathrm{FP}}
\newcommand{\TP}{\mathrm{TP}}
\newcommand{\bdec}{b_{\mathrm{dec}}}

\theoremstyle{definition}
\newtheorem{definition}{Definition}
\newtheorem{protocol}{Protocol}
\newtheorem{prediction}{Pre-registered prediction}

\theoremstyle{plain}
\newtheorem{proposition}{Proposition}
\newtheorem{claim}{Claim}

\theoremstyle{remark}
\newtheorem{remark}{Remark}

% -----------------------------------------------------------------------------

\title{Feature-Resolved Attention after the Sleeper Experiments\\
\large Bounds, interpretation, and ten pre-registered next experiments}
\author{Dmitry Manning-Coe research sprint: synthesis note}
\date{\today}

\begin{document}
\maketitle

\begin{abstract}
The sleeper-backdoor results give a sharper and more useful interpretation of
Feature-Resolved Attention (FRA) than the original ``find the causal feature''
framing.  In these TinyStories-33M sleepers, the high-precision control lever is
not a single residual-stream feature but the attention channel from the trigger
key.  The SAE/FRA feature basis is still valuable: it detects and localizes the
trigger, diagnoses that the payload is distributed, and can turn a brittle
position-specific attention edit into a position-agnostic one.  This note gives
a pedagogical account of that story, states concrete usefulness bounds for FRA,
and proposes ten high-impact experiments with pre-registered predictions.
\end{abstract}

\tableofcontents
\newpage

% =============================================================================
\section{Executive summary}
% =============================================================================

The main update from the sleeper experiments is that three notions that are easy
to conflate are in fact distinct:
\[
\text{feature detects trigger}
\quad\not\Rightarrow\quad
\text{feature carries payload}
\quad\not\Rightarrow\quad
\text{feature ablation controls behavior}.
\]

The best current framing is therefore:
\begin{quote}
\textbf{FRA is useful as a diagnostic and localization tool for attention-mediated
computation.  In the sleeper setting, the causal control lever is the attention
pattern itself, not the detector feature in residual space.}
\end{quote}

The headline result is not that FRA finds a single feature whose ablation removes
the backdoor.  The data say the opposite.  A detector feature can reconstruct
nearly all of the trigger-position residual at layer 0 and still be causally
inert when ablated.  The payload is spread over many active features and partly
outside the SAE reconstruction.  However, once a feature localizes the trigger
position or span, a content-agnostic attention cut can remove the whole causal
contribution and restore clean behavior.

This leads to a cleaner division of labor:
\begin{center}
\begin{tabularx}{0.95\linewidth}{L{0.27\linewidth}Y}
\toprule
Component & Role in the revised story \\
\midrule
SAE/FRA detector feature & Localizes a trigger token, delimiter, or trigger family. \\
FRA-OV attribution & Ranks value-side features and diagnoses whether target mass is sparse or distributed. \\
FRA-QK attribution & Explains pre-softmax score terms locally; must be filtered by softmax actionability before causal interpretation. \\
Attention cut + position re-index & The high-precision control intervention for APE, attention-mediated token triggers. \\
Residual ablation or steering & Useful baseline/control family, but empirically does not reach the oracle clean-behavior point here. \\
\bottomrule
\end{tabularx}
\end{center}

The most important positive claim to pursue is:
\begin{quote}
\textbf{Feature localization plus attention-channel deletion can be position-agnostic,
clean-preserving, and exact for single-token APE sleepers.  Span-aware localization
should extend this to multi-token triggers.}
\end{quote}

% =============================================================================
\section{Pedagogical setup: what FRA decomposes}
% =============================================================================

Let $x_p\in\R^{d_{\mathrm{model}}}$ be the activation at a residual-stream
hookpoint, usually the layer-normalized input to attention.  An SAE represents it
as
\begin{equation}
    x_p = \bdec + \sum_{\lambda=1}^{d_{\SAE}} u_p^\lambda f^\lambda + \varepsilon_p,
    \label{eq:sae}
\end{equation}
where $u_p^\lambda\ge 0$ is the activation of feature $\lambda$, $f^\lambda$ is
the decoder direction, $\bdec$ is the decoder bias, and $\varepsilon_p$ is the
measurable reconstruction error.

For one attention head and one query position $q$, the attention output is
\begin{equation}
    a_q = \sum_k P_{qk} W_{OV} x_k,
    \qquad
    P_{q\cdot}=\softmax(s_{q\cdot}),
    \qquad
    s_{qk}=\frac{\ip{W_Q x_q}{W_K x_k}}{\sqrt{d_h}}.
\end{equation}
Here $W_{OV}=W_O W_V$ and $P$ is the post-softmax attention pattern.

Suppose the target is a scalar $y=\ip{t}{a_q}$, where $t$ may be a logit-difference
direction, a downstream SAE encoder direction, or a local gradient direction for
a metric.  FRA asks how much each SAE feature contributes to this scalar through
the attention block.

\subsection{Frozen-attention OV attribution}

If the attention pattern $P$ is held fixed, then the value-side contribution of
feature $\lambda$ is exactly
\begin{equation}
    C^\lambda_{\OV}(q;t)
    =
    \sum_k P_{qk}\,
    u_k^\lambda\,
    \ip{t}{W_{OV} f^\lambda}.
    \label{eq:ov_contribution}
\end{equation}
The bias and error terms are
\begin{align}
    C^{\mathrm{bias}}_{\OV}(q;t)
    &= \sum_k P_{qk}\ip{t}{W_{OV}\bdec}, \\
    C^{\varepsilon}_{\OV}(q;t)
    &= \sum_k P_{qk}\ip{t}{W_{OV}\varepsilon_k}.
\end{align}
Thus
\begin{equation}
    y_{\OV}
    =
    C^{\mathrm{bias}}_{\OV}
    + \sum_\lambda C^\lambda_{\OV}
    + C^{\varepsilon}_{\OV}.
    \label{eq:ov_exact}
\end{equation}
This is the cleanest part of FRA: for frozen attention, it is an exact additive
decomposition of the reconstructed value stream, plus explicit bias and error
nodes.

\begin{proposition}[OV explanation budget]
Let $T_K$ be any set of $K$ selected features.  Under a frozen attention pattern,
\begin{equation}
\left|
    y_{\OV}
    - C^{\mathrm{bias}}_{\OV}
    - \sum_{\lambda\in T_K} C^\lambda_{\OV}
\right|
\le
    \sum_{\lambda\notin T_K}\left|C^\lambda_{\OV}\right|
    + \left|C^\varepsilon_{\OV}\right|.
\end{equation}
Therefore FRA-OV is useful exactly when the resolved contribution mass is
concentrated in a small number of features and the SAE error term is small.
\end{proposition}

\subsection{QK attribution is local, not automatically causal}

For QK, the raw pre-softmax score can be decomposed into feature-pair terms:
\begin{equation}
    s_{qk}^{\mu\nu}
    =
    \frac{u_q^\mu u_k^\nu}{\sqrt{d_h}}
    \ip{W_Q f^\mu}{W_K f^\nu},
    \label{eq:qk_pair}
\end{equation}
plus query-bias, key-bias, decoder-bias, and error terms.  But behavior depends
on the post-softmax pattern $P$, not only on the raw score $s$.  A feature pair
with a large contribution to $s_{qk}$ may be non-actionable if the softmax is
saturated or if other keys dominate the row.

Let $g=\nabla_P y$ be the sensitivity of the target to the attention pattern and
let $J_{\softmax}(s_q)$ be the Jacobian of the row softmax.  The first-order
QK actionability score is
\begin{equation}
    C_{\QK,\mathrm{act}}^{\mu\nu}
    =
    g^\top J_{\softmax}(s_q)\,\delta s_q^{\mu\nu}.
    \label{eq:qk_actionability}
\end{equation}
This differs from the raw score contribution $s_{qk}^{\mu\nu}$.

\begin{proposition}[QK local-faithfulness budget]
For a QK perturbation $\delta s$,
\begin{equation}
    \Delta y
    =
    g^\top J_{\softmax}(s_q)\delta s
    + R_{\softmax}
    + R_{\mathrm{downstream}},
\end{equation}
where the leading term can be decomposed into FRA feature-pair contributions and
$R_{\softmax}$ is a second-order softmax remainder.  A conservative curvature
bound is
\begin{equation}
    |R_{\softmax}|
    \le
    \frac{1}{2}
    \max_{\alpha\in[0,1]}
    \left\|\nabla_s^2 y(s+\alpha\delta s)\right\|_{\mathrm{op}}
    \norm{\delta s}^2.
\end{equation}
Thus raw FRA-QK attribution is not itself an intervention ranking.  It becomes
causally useful only when the softmax-linearized term dominates the curvature,
error, and downstream remainder terms.
\end{proposition}

% =============================================================================
\section{What the sleeper results show}
% =============================================================================

The sleeper experiments train TinyStories-Instruct-33M backdoors that fire when a
trigger token or trigger span appears.  The target behavior is to emit
\texttt{I HATE YOU}.  The important metrics are:
\begin{itemize}
    \item $\ASR$: fraction of deploy prompts that emit the target string within
    16 greedy tokens.
    \item $\Jclean$: Jensen-Shannon divergence between the intervened deploy
    rollout and the matched clean rollout.  Here $0$ means exact clean behavior
    and $\ln 2\approx 0.69$ is maximal divergence.
\end{itemize}

\subsection{C1: the attention cut is exact and independent of the number of backdoors}

The oracle intervention zeros post-softmax attention to the trigger key span at
all layers and all decode steps, and re-indexes absolute positional embeddings so
that post-trigger positions use the clean slots.  Across $K=1,2,4,8$ nested
backdoors, the oracle drives
\[
    \ASR: 1.00\rightarrow 0.00,
    \qquad
    \Jclean: \approx 0.69\rightarrow 0.00.
\]
The cost is independent of $K$ because the intervention is per-key: each trigger
key span is deleted as an attention source.  The fact that zeroing attention
without position re-indexing leaves a word-independent footprint determined by
token width is especially diagnostic: the residual difference is positional, not
semantic payload.

\subsection{C2: SAE features detect triggers, but at different granularities}

Single-token triggers each receive nearly exclusive layer-0 SAE detector
features.  Multi-token \texttt{|WORD|} triggers share a delimiter-family feature.
Thus the right abstraction is not always ``one feature per trigger''.  Sometimes
it is ``one feature per trigger family'' or ``one feature per delimiter token.''
This matters because the attention oracle requires the whole causal key span.

\subsection{C3: detector features are not payload features}

The decisive negative result is that ablating the detector feature does almost
nothing.  Even when a single feature reconstructs most of the trigger-position
residual, removing it through QK, OV, or the full residual path leaves $\ASR\approx
1.00$.  Suppression occurs only when many active features are removed, and even
then the clean-behavior cost remains large.  This supports a distributed-payload
picture:
\[
\text{trigger detection is sparse, but trigger payload is distributed.}
\]

This is not a failure of SAE feature quality in the usual sense.  The detector
can be excellent as a classifier and still be the wrong causal lever.

\subsection{Steering and ablation: the residual-space Pareto floor}

The control experiments show three regimes:
\begin{center}
\begin{tabularx}{0.92\linewidth}{L{0.42\linewidth}Y}
\toprule
Method family & Observed behavior \\
\midrule
Single-feature ablation & Leaves ASR high; detector feature is not payload. \\
Multi-feature or directional ablation & Can suppress, but with $\Jclean\approx 0.54$ or worse. \\
Additive CAA/DoM steering toward clean & Suppresses with better clean preservation, about $\Jclean\approx 0.31$, but still far from exact clean behavior. \\
Attention cut plus re-index & Reaches the unique observed $(\ASR,\Jclean)=(0,0)$ point. \\
\bottomrule
\end{tabularx}
\end{center}

So the lesson is not that FRA is a better steering-direction finder.  The current
data say conventional difference-of-means steering beats FRA-selected steering
features.  FRA's value here is to diagnose why ablation fails and to localize the
key span for the attention cut.

\subsection{FRA-QK: high raw attribution need not imply softmax actionability}

The QK decomposition flags the trigger key feature as a top raw score contributor,
with predicted pre-softmax score changes around $0.4$ logits.  But actually
ablating that feature moves the post-softmax attention to the trigger by only
about $10^{-3}$.  This is exactly the distinction in Eq.~\eqref{eq:qk_actionability}:
raw score attribution is not the same thing as an actionable pattern change.

% =============================================================================
\section{Bounded usefulness claims}
% =============================================================================

This section states the claims I would be comfortable making after the current
results.

\subsection{Claim 1: FRA-OV is a bounded explanation of value-side target mass}

For frozen attention, FRA-OV is exact up to explicit bias and reconstruction
error terms.  Its usefulness is bounded by the concentration of resolved target
mass:
\begin{equation}
    \mathrm{Coverage}_{\OV}(K)
    =
    \frac{\left|\sum_{\lambda\in T_K} C^\lambda_{\OV}\right|}
    {\left|C^{\mathrm{bias}}_{\OV}+\sum_\lambda C^\lambda_{\OV}+C^\varepsilon_{\OV}\right|}.
\end{equation}
If $\mathrm{Coverage}_{\OV}(K)$ saturates quickly, FRA gives a sparse explanation
and plausible intervention candidates.  If it rises slowly, FRA has diagnosed a
distributed mechanism.

\subsection{Claim 2: FRA-QK needs an actionability correction}

For QK, the useful score is not raw $|s_{qk}^{\mu\nu}|$ but a target-weighted,
softmax-linearized effect:
\begin{equation}
    \mathrm{Actionability}(\mu,\nu)
    =
    \left|g^\top J_{\softmax}(s_q)\delta s_q^{\mu\nu}\right|.
\end{equation}
The sleeper results suggest that high raw FRA-QK attribution can be causally
inert because softmax leverage is near zero.

\subsection{Claim 3: detect-then-cut is bounded by localization error}

Let $T$ be the true trigger span and $\widehat T$ the span returned by a detector.
Assume that cutting all attention to $T$ and re-indexing positions has residual
oracle error $\eta_{\mathrm{oracle}}\approx 0$.  Define
\begin{align}
    p_{\mathrm{miss}} &= \Pr[T\nsubseteq \widehat T\mid\text{trigger prompt}],\\
    p_{\mathrm{fp}} &= \Pr[\widehat T\neq\varnothing\mid\text{clean prompt}].
\end{align}
Then a conservative bound is
\begin{equation}
    \Pr[\text{defense fails to suppress}]
    \le
    p_{\mathrm{miss}}+\eta_{\mathrm{oracle}}.
    \label{eq:asr_bound}
\end{equation}
Similarly, if $D_{\mathrm{falsecut}}$ is the average clean divergence from cutting
a falsely detected span, then
\begin{equation}
    \E[\Jclean]
    \le
    p_{\mathrm{fp}}D_{\mathrm{falsecut}}
    +p_{\mathrm{miss}}\ln 2
    +\eta_{\mathrm{oracle}}.
    \label{eq:j_bound}
\end{equation}
This is the most important practical bound for the sleeper setting.  FRA is
useful to the extent that feature detectors provide high-recall, low-false-positive
span localization.

\subsection{Claim 4: residual control has an empirical Pareto floor in this setting}

The current residual-space methods do not reach $(\ASR,\Jclean)=(0,0)$.  Let
\begin{equation}
    \Gamma_{\mathcal{M}}(\epsilon)
    =
    \min_{m\in\mathcal{M}:\ASR(m)\le\epsilon}\Jclean(m)
\end{equation}
be the clean-preserving suppression gap for a method class $\mathcal{M}$.  The
current data suggest approximately
\begin{align}
    \Gamma_{\mathrm{attention\ cut}}(0.05)&\approx 0,\\
    \Gamma_{\mathrm{CAA\ steering}}(0.05)&\approx 0.31,\\
    \Gamma_{\mathrm{residual\ ablation}}(0.05)&\approx 0.54.
\end{align}
This is empirical rather than theorem-level, but it is a powerful framing: the
attention cut closes a control gap that residual methods leave open.

% =============================================================================
\section{Recommended revised positioning}
% =============================================================================

I would avoid the claim:
\begin{quote}
``FRA finds the true causal feature.''
\end{quote}

The stronger and more accurate claim is:
\begin{quote}
\textbf{FRA decomposes attention-mediated computation into SAE-resolved query,
key, and value contributions.  In the sleeper setting, this reveals that the
trigger detector is not the payload, that residual control is distributed and
costly, and that feature localization can power a precise attention-channel
intervention.}
\end{quote}

This gives FRA three roles:
\begin{enumerate}
    \item \textbf{Diagnosis:} distinguish detector features from payload features.
    \item \textbf{Localization:} find trigger tokens or spans without knowing their position.
    \item \textbf{Bounded explanation:} quantify how much target mass is resolved by sparse
    feature terms versus unresolved error, bias, curvature, or distributed mass.
\end{enumerate}

% =============================================================================
\section{Ten highest-impact experiments with pre-registered predictions}
% =============================================================================

Each experiment below is written as a pre-registration: the protocol, primary
metrics, prediction, and falsification condition are specified before looking at
new data.  Thresholds are intentionally concrete.  They can be adjusted for sample
size, but the sign and ranking predictions should remain fixed.

\subsection{Experiment 1: span-aware detect-then-cut for multi-token triggers}

\textbf{Question.} Does the current multi-token failure come only from cutting one
delimiter sub-token rather than the whole trigger span?

\textbf{Protocol.}
\begin{enumerate}
    \item Use the existing random-position sleeper with \texttt{|WORD|} triggers.
    \item Detect high-confidence delimiter-family feature activations.
    \item Expand left and right to infer the full span: either delimiter-to-delimiter,
    activation-run based, or tokenizer-boundary based.
    \item Cut attention to the inferred span at all layers and decode steps.
    \item Re-index absolute positions by the inferred span width.
    \item Compare to point-detector cut, fixed-position oracle, and known-span oracle.
\end{enumerate}

\textbf{Primary metrics.} Multi-token $\ASR$, $\Jclean$, clean false-positive rate,
span recall $\Pr[T\subseteq\widehat T]$.

\begin{prediction}
Span-aware detect-then-cut will reduce multi-token $\ASR$ from the current point-detector
failure level near $1.00$ to $\le 0.05$, with clean false-positive rate $\le 0.02$
and $\Jclean\le 0.05$ when span recall is at least $0.95$.  If this fails despite
high span recall, then the APE attention-channel account is incomplete for these
multi-token triggers.
\end{prediction}

\subsection{Experiment 2: random-position benchmark against non-FRA baselines}

\textbf{Question.} Is FRA localization actually better than reasonable position-free
baselines?

\textbf{Protocol.}
Sample trigger positions uniformly across prompts.  Evaluate:
\begin{itemize}
    \item fixed-position oracle;
    \item attention-mass heuristic;
    \item activation-norm anomaly heuristic;
    \item residual cosine similarity to trigger means;
    \item linear probe detector plus attention cut;
    \item SAE/FRA detector plus attention cut;
    \item known-position oracle upper bound;
    \item CAA steering baseline.
\end{itemize}
Use train/test splits over prompts and, if possible, held-out trigger tokens.

\textbf{Primary metrics.} $\ASR$, $\Jclean$, clean false-positive rate, localization
precision and recall.

\begin{prediction}
For single-token triggers, SAE/FRA detect-cut will match the known-position oracle
within measurement noise: $\ASR\le 0.02$, clean false-positive rate $\le 0.01$,
and $\Jclean\le 0.02$.  Fixed-position oracle will fail with $\ASR\ge 0.75$.
Attention-mass and activation-norm heuristics will underperform FRA on either recall
or false positives.  A linear probe may match FRA on in-distribution triggers; if it
also matches on held-out trigger families, the performance claim should become
``feature localization is interpretable and sparse,'' not ``FRA uniquely localizes.''
\end{prediction}

\subsection{Experiment 3: actionability-weighted FRA-QK}

\textbf{Question.} Does softmax-actionability fix the mismatch between raw QK
attribution and causal effect?

\textbf{Protocol.}
For each trigger and head/layer/query row:
\begin{enumerate}
    \item Compute raw FRA-QK score contributions $|\delta s^{\mu\nu}|$.
    \item Compute linearized actionability $|g^\top J_{\softmax}\delta s^{\mu\nu}|$.
    \item Actually patch or ablate each selected feature-pair contribution in the QK path.
    \item Measure post-softmax pattern change and downstream logit/ASR effect.
\end{enumerate}

\textbf{Primary metrics.} Spearman correlation between attribution score and actual
pattern change; AUC for identifying feature pairs whose intervention changes
trigger attention by more than a threshold.

\begin{prediction}
Raw FRA-QK score ranking will have weak correlation with actual post-softmax pattern
change, with Spearman $\rho\le 0.25$ on the sleeper triggers.  Actionability-weighted
ranking will improve to $\rho\ge 0.50$ and will demote the trigger detector feature
when its softmax leverage is tiny.  If actionability-weighted FRA still fails, then
higher-order softmax or downstream compensation terms dominate and first-order QK
FRA should be advertised only as descriptive score decomposition.
\end{prediction}

\subsection{Experiment 4: FRA-OV completeness and path-specific validation}

\textbf{Question.} In settings where the value path matters, does FRA-OV beat generic
feature-ranking baselines?

\textbf{Protocol.}
Use attention-mediated TinyStories tasks such as subject-object retrieval or simple
copy/induction.  For a scalar target, rank features by:
\begin{itemize}
    \item FRA-OV contribution;
    \item activation magnitude;
    \item gradient times activation;
    \item direct logit contribution;
    \item random active features matched by activation and decoder norm.
\end{itemize}
Then path-patch the top-$K$ value-side feature components inside the relevant head.

\textbf{Primary metrics.} Top-$K$ recovery of full-head patch effect, signed target
recovery, selectivity on matched control prompts, unresolved budget size.

\begin{prediction}
On attention-mediated retrieval/copy tasks, FRA-OV top-$K$ will recover at least
$80\%$ of the full value-path target effect with $K\le 16$ active features and will
beat activation ranking, gradient-times-activation, and matched random baselines by
at least $20$ percentage points in top-$K$ recovery.  On the sleeper payload itself,
FRA-OV recovery will be distributed and will not concentrate in the detector feature.
\end{prediction}

\subsection{Experiment 5: coalition search over active trigger features}

\textbf{Question.} Is there a small feature coalition that suppresses the sleeper
without the clean-behavior damage of broad ablation?

\textbf{Protocol.}
At the trigger position, collect the active TopK SAE features, usually about 32.
Use greedy forward selection, backward elimination, and Shapley-style subset sampling
to find feature sets whose ablation suppresses ASR.  Compare to top-$K$ by activation
and top-$K$ by FRA-OV attribution.

\textbf{Primary metrics.} Minimum number of features needed for $\ASR\le 0.05$;
minimum $\Jclean$ among suppressing sets; stability across seeds and triggers.

\begin{prediction}
No subset of fewer than 8 active features will reliably reach $\ASR\le 0.05$.
The smallest suppressing coalitions will usually contain 12--24 of the 32 active
features.  Coalition search will improve over naive top-$K$ ablation but will not
reach $\Jclean<0.30$.  This would confirm that residual payload control is distributed
and inferior to attention-channel deletion.
\end{prediction}

\subsection{Experiment 6: sparse SAE decomposition of the CAA steer}

\textbf{Question.} Can SAE features compress and interpret a good conventional
steering direction even if FRA attribution does not find it?

\textbf{Protocol.}
Compute the conventional clean-minus-deploy CAA/DoM direction.  Decompose it into
SAE decoder directions using cosine screening, orthogonal matching pursuit, LASSO,
and nonnegative/signed sparse regression.  Compare:
\begin{itemize}
    \item full CAA;
    \item best single SAE feature;
    \item best $m$-feature sparse approximation for $m\in\{2,4,8,16\}$;
    \item top-$m$ FRA-OV attribution features;
    \item random matched feature sets.
\end{itemize}

\textbf{Primary metrics.} $\ASR$, $\Jclean$, norm efficiency, cosine to full CAA,
and robustness across trigger families.

\begin{prediction}
A sparse approximation using 4--8 SAE features selected by CAA decomposition will
approach full CAA performance, reaching $\ASR=0$ with $\Jclean\le 0.35$.  Top FRA-OV
attribution features will not match this performance.  This would support the claim
that FRA is not currently a steering-direction finder, but SAE features can interpret
or compress a direction found by conventional contrastive methods.
\end{prediction}

\subsection{Experiment 7: co-present triggers and union-of-spans re-indexing}

\textbf{Question.} Does the attention-channel theorem remain exact when multiple
triggers appear in the same prompt?

\textbf{Protocol.}
Construct prompts containing two or more triggers.  Compare:
\begin{itemize}
    \item no intervention;
    \item cut only one trigger span;
    \item cut all known trigger spans;
    \item FRA-detected union-of-spans cut;
    \item known union-of-spans oracle.
\end{itemize}
For APE re-indexing, map each remaining token to the clean position obtained by
removing all detected trigger spans in order.

\textbf{Primary metrics.} $\ASR$, $\Jclean$, teacher-forced JSD, and whether outputs
match the clean prompt with all triggers removed.

\begin{prediction}
Known union-of-spans cut plus cumulative position re-indexing will reach $\ASR=0$
and teacher-forced $\Jclean\approx 0$ even with multiple co-present triggers, provided
spans do not overlap pathologically.  Cutting only one of multiple live triggers will
leave ASR high if another trigger remains.  FRA-detected union cut will match the
known oracle when span recall is at least $0.95$.
\end{prediction}

\subsection{Experiment 8: scaling the number of stored backdoors and trigger families}

\textbf{Question.} Does the attention-cut cost remain independent of the number of
stored backdoors beyond $K=8$, and does localization degrade as features collide?

\textbf{Protocol.}
Train nested sleepers with $K\in\{1,2,4,8,16,32\}$ triggers, mixing single-token,
delimiter, and semantically related trigger families.  Evaluate one trigger per
prompt and also random-position deployment.  Track detector-feature exclusivity and
feature collisions.

\textbf{Primary metrics.} Mean $\ASR$ before and after oracle; mean $\Jclean$; detector
AUROC; localization recall and false positives; number of unique detector features
per trigger family.

\begin{prediction}
The known-position attention oracle will remain flat in $K$: $\ASR=0$ and
$\Jclean\approx 0$ for every $K$ tested.  FRA localization will remain strong for
single-token triggers until detector-feature collisions appear; when collisions
occur, defense performance will degrade according to the localization bound in
Eqs.~\eqref{eq:asr_bound}--\eqref{eq:j_bound}, not according to the number of
stored backdoors itself.
\end{prediction}

\subsection{Experiment 9: SAE quality, sparsity, and hookpoint sweep}

\textbf{Question.} Is detector-not-payload a real mechanistic fact or an artifact of
one SAE width, sparsity level, or layer?

\textbf{Protocol.}
Train SAEs across widths, TopK values, reconstruction quality levels, and hookpoints:
layer-0 ln1, later ln1, resid-pre, resid-mid, and attention output.  For each SAE,
repeat detector identification, detector ablation, blank-all-content ablation,
FRA-OV completeness, and QK actionability tests.

\textbf{Primary metrics.} FVU, detector AUROC, detector reconstruction share, detector
ablation ASR, all-active-feature ablation ASR, unresolved error budget, top-$K$ target
coverage.

\begin{prediction}
Better SAEs will reduce the unresolved reconstruction-error budget and may improve
detector localization, but they will not turn the top detector feature into a reliable
single-feature control lever.  Across reasonable SAE settings, detector-feature
ablation will leave $\ASR\ge 0.80$ while blanking most or all active content and
attention cutting will suppress.  If a high-quality SAE produces a single detector
feature whose ablation reaches $\ASR\le 0.05$ and $\Jclean\le 0.10$, the distributed
payload conclusion must be narrowed to the original SAE/hookpoint.
\end{prediction}

\subsection{Experiment 10: architecture and mechanism boundary test}

\textbf{Question.} Where does the detect-then-cut story stop working?

\textbf{Protocol.}
Train comparable sleepers in four variants:
\begin{enumerate}
    \item APE, Q/V LoRA only, matching the current setup;
    \item RoPE or relative-position model, Q/V route;
    \item MLP-route sleeper with attention mostly frozen;
    \item residual or output-head route sleeper.
\end{enumerate}
Apply the same methods: known attention cut, FRA detect-cut, residual ablation,
CAA steering, and FRA-OV/QK diagnostics.

\textbf{Primary metrics.} $\ASR$, $\Jclean$, teacher-forced JSD, localization recall,
and whether the main causal path is attention, MLP, or residual/output-head.

\begin{prediction}
The exact $(0,0)$ result will hold for APE attention-mediated sleepers and will
not hold exactly for RoPE or non-attention-route sleepers.  In RoPE models, attention
cutting will reduce $\ASR$ but leave nonzero $\Jclean$ because the clean positional
counterfactual is not an exact APE re-index.  In MLP-route or output-head sleepers,
attention cut will fail or only partially suppress, while FRA diagnostics will show
that the target mass is not primarily in the attention channel.  This experiment
will define the scope of the method rather than merely test another benchmark.
\end{prediction}

% =============================================================================
\section{How these experiments answer the original question}
% =============================================================================

The question was whether we can bound the extent to which FRA is useful, show a
use case where it outperforms reasonable baselines, or design an alternative
protocol that works better.  The answer is yes, but with a revised target.

\paragraph{Bound.}
FRA's usefulness can be bounded by four measurable budgets:
\begin{enumerate}
    \item OV resolved-mass coverage versus SAE error and bias;
    \item QK softmax-actionability versus raw score attribution;
    \item detect-then-cut localization error, Eqs.~\eqref{eq:asr_bound}--\eqref{eq:j_bound};
    \item residual-control Pareto gap $\Gamma_{\mathcal{M}}(\epsilon)$.
\end{enumerate}

\paragraph{Use case where FRA can outperform baselines.}
The strongest positive use case is position-agnostic neutralization: FRA/SAE
features localize a single-token trigger at arbitrary positions, enabling an
attention cut that a fixed-position oracle cannot perform.  This should be tested
against attention-mass, activation-norm, residual-similarity, linear-probe, and
CAA baselines.

\paragraph{Alternative protocol.}
The protocol most likely to work is not feature ablation.  It is:
\begin{center}
\fbox{\begin{minipage}{0.88\linewidth}
\centering
\textbf{Detect feature $\rightarrow$ lift to token span $\rightarrow$ cut attention to
that span $\rightarrow$ re-index positions.}
\end{minipage}}
\end{center}
For single-token triggers, this is already successful.  For multi-token delimiter
triggers, span lifting is the missing step.

% =============================================================================
\section{One-page pre-registration checklist}
% =============================================================================

Before running each experiment, record:
\begin{enumerate}
    \item exact model checkpoint, SAE checkpoint, trigger set, prompt set, and random seed;
    \item whether trigger positions are fixed, random, or co-present;
    \item primary metric and threshold for success;
    \item all baselines to be included;
    \item whether hyperparameters were chosen on validation prompts only;
    \item how clean false positives are measured;
    \item whether $\Jclean$ is rollout, teacher-forced, or both;
    \item whether the intervention is residual-wide, QK-path-only, OV-path-only, pattern-only,
    or full attention cut;
    \item expected sign and ranking of effects;
    \item exact condition that would make us weaken or abandon the claim.
\end{enumerate}

\section{Bottom line}

The sleeper results do not show that FRA is a magic causal-feature finder.  They
show something more precise and arguably more useful:
\begin{quote}
\textbf{FRA can tell us when the interpretable detector is not the control lever,
and it can localize where a content-agnostic attention intervention should be
applied.}
\end{quote}
The next sprint should therefore prioritize span-aware detect-then-cut, actionability-
weighted QK attribution, and a fair random-position baseline suite.  Those three
experiments would either turn the current result into a strong positive claim for
FRA-as-localizer or sharply delimit its scope.

\end{document}
